\subsection*{先分清同一笔付款与覆盖新状态}

W1 在第 7 代发出 \code{pay-C1}，W2 在第 8 代继续恢复 C1，两次调用的目标相同：完成原来受理的那一笔全额付款。平台按原范围、原键和原输入合流两次到达，才能共同找到 X1。领取换代改变了调度者，付款授权继续绑定原 C1。

覆盖写面对另一种问题。W2 已把报告更新为新阶段，W1 先前准备的旧阶段随后到达，资源需要判断哪次领取已经过时。接收者为此比较代次，并把比较与实际写入连在一起。本节先推导这项资源保护，再检查同一代次内不同动作的重复和乱序。

\subsection*{本地查过资格以后可能发生什么}

前面的 P/C/T、观察与支付结果都在购买库中，由 C→T 的本地事务和相应入口保护。现在另设一个展示进度的旁例：运维界面从报告服务读取 \code{report:C1}，W1、W2 把 T1 的核查进度写给这个服务。报告服务拥有独立存储，方便界面单独读取和更新展示内容，因此这次覆盖写跨出了购买库的事务边界。支付成功与拒绝继续由原 C 的证据决定。第十章向购买库回报履约结果的待办是另一项业务对象。

先取 \code{report:C1} 只由 T1 管理，已接受代次为 7，展示内容 a 表示核查已开始。W1 在本地读到第 7 代有效，准备再次写入 a。随后 W1 暂停，W2 接管第 8 代，取得核查进展，准备把展示更新为 b。W2 的 b 先到，W1 恢复后又把 a 送达。独立存储只接收内容，尚无本地领取发生变化的资料，于是把 b 覆盖为 a。把本地资格检查挪到发送前最后一行，暂停也可能发生在检查与发送之间，请求还可能在网络中延迟。

若希望接收端拒绝这种旧写，请求应携带能够比较领取先后的凭据。本章把它称为 fencing token，即让资源排除旧领取写入的代次标记。对只由 T1 管理的资源 \code{report:C1}，可以使用带作用域的 \code{(T1, lease\_epoch)}。资源同时核验发送者具有该作用域的权限，避免任意客户伪造更大数字。

资源保存 \code{high\_water}，表示它已经接受的最大代次。这个字段由资源自己的受保护写事务更新，每次写入前读取。若传入代次小于水位，拒绝写。大于或等于水位，按业务规则处理，并把水位推进到二者的较大值。资源名称、T1 的映射和水位都持久保存，重启后继续承担拒绝依据。

\subsection*{高水位从什么时候开始起作用}

先取 \code{report:C1} 水位为 7，内容为旧工作阶段。租约服务已经授予 W2 第 8 代，但资源尚未收到任何第 8 代请求。W1 的第 7 代先到，比较 7 与当前水位 7，可以通过。W2 的第 8 代随后到，资源接受并推进水位为 8。再到达的第 7 代才被拒绝。

\CEightFigure{ch08-original06}{最近已观察水位在资源接受更高代次以后排除旧写}{fig:ch08-watermark}

这一规则保证：资源一旦接受第 8 代，以后就拒绝较小代次的写入。它的保护起点位于资源接受新代次之时。租约服务发出 8，与资源观察到 8 之间可以存在间隔，只保存最近已观察水位的接收端，在间隔中继续根据水位 7 作判断。

有的业务要求 W2 接管完成时，资源就已经排除旧写。可以把接管定义为两个阶段：先取得候选代次 8，再在资源原子安装水位 8，取得确认后才向 W2 报告可使用资源。若确认丢失，W2 读取资源：水位等于 8，可以恢复这次安装；水位小于 8，沿原候选重试安装；水位已经大于 8，则结束旧候选。安装请求本身也经过权限和代次检查。正常确认只说明这次安装已经成立，后续也可能出现更高代次，每次业务写入都携带代次，由资源在原子边界核验。

\CEightFigure{ch08-install}{把资源安装确认纳入接管完成，才在报告可用时排除旧代次}{fig:ch08-install}

这样定义的可用时点包含资源安装，旧写若先于安装完成，可以在资源顺序中排在它之前。若进一步要求从租约服务分配 8 的那个瞬间起排除所有第 7 代写入，分配与资源准入就必须共享仲裁边界，或由资源采用与分配服务协调的校验机制。单独远程查询一次最新代次，随后无条件写入，会重现检查后的间隔。

\subsection*{检查与写入在哪一处相连}

资源即使保存水位，也可能错误地把读取和修改拆开。W1 读到水位 7，通过检查。W2 接受 8 并提交新内容。W1 随后按先前检查结果直接写旧内容。最终水位可以为 8，内容却被旧请求覆盖。只观察最大数字，会漏掉这条历史。

资源把检查、水位推进和业务写入放在同一个原子决定中，才能覆盖内容更新。若资源本身是数据库，可以锁住资源行，比较传入代次，随后更新水位和内容，共同提交。另一种可用表达是包含条件的 UPDATE，前提是数据库确实在并发冲突后按当前行执行条件，并且相关业务写入都由这项原子条件覆盖~\cite{pg}。

\begin{BookCode}
UPDATE resource_state
SET high_water = :epoch,
    value = :new_value
WHERE resource_id = 'report:C1'
  AND high_water <= :epoch;
\end{BookCode}

这里 \code{:epoch} 已经过作用域与来源核验，所有更新都走这个入口，水位初值也由受控创建确定。对单行内容，影响一行则提交水位和内容，影响零行表示代次过旧或资源缺失，服务读取具体状态后分别返回旧写拒绝或关联异常。数据库异常则回滚。提交确认丢失时，读取资源水位及对应业务操作记录判断原提交，重复动作的识别规则见下文。若业务效果涉及多张表，则在同一资源事务中锁行检查，再完成整组写入，所有效果共同受这次检查约束。

\CEightFigure{ch08-original07}{高水位检查与内容提交必须共享资源端的原子边界}{fig:ch08-atomic-resource}

可以逐次说明水位规则为什么成立。资源初始水位为 0，受控创建尚无旧持有者。每次接受写入都要求传入 \(e\geq h\)，其中 \(h\) 是当前水位，提交后水位变为 \(e\)，所以水位保持不下降。某次已接受 \(e=8\) 后，传入 7 的检查必然失败。检查与业务提交共同仲裁，使这个结论覆盖内容更新。资源重启保留水位，才能把同一推理延续到重启后的迟到请求。

\subsection*{相同 token 怎样与业务身份配合}

代次水位已经说明怎样排除旧领取。接着在同一个第 8 代内部逐步增加问题：先让同一动作重复，再让两个不同进度倒序到达。这两条历史分别检查业务身份与业务版本，所用资源继续是前面的 \code{report:C1}。

同一工作者可能在第 8 代内多次发送相同业务动作。资源按 \(e\geq h\) 接受两次代次 8，若动作都是 \code{counter=counter+1}，计数便增加两次。要使这项动作重复到达时只产生一次效果，资源在通过代次检查后，再用业务操作身份查找已有提交，并为首次动作共同保存效果与原结果。

把水位条件改成严格大于，会损害正常工作：第 8 代第一次修改后，后续不同的合法动作都将被拒绝。需要多次动作的资源，可以在一个事务中同时检查代次与业务 \code{operation\_id}。以 \code{operation-C1} 为例，重复尝试恢复原提交，同一代次内另一项明确授权的动作取得不同身份。

再让 W2 在第 8 代先发送 a、后发送 b，两个动作各有自己的 \code{operation\_id}。b 先到，a 迟到，两次代次检查都满足 \(8\geq 8\)，两个操作号也都第一次出现。只做上述检查，展示就会从 b 退回 a。代次比较排除旧领取，操作号识别重复动作，同一代次内不同进度的业务先后还要由进度规则判断。

对这个展示对象，可以把每份内容定义为完整进度快照，并附上 \code{progress\_version}，标识 T1 这项核查的业务进度版本。为产生这样的版本，购买服务作为报告生产方，在购买库中另设一份以 T1 为键的核查进度记录，初始版本为 0。工作者推进核查阶段时，沿 C→T 的锁定顺序核验当前领取资格，再锁定这份进度记录，核对合法前驱，把进度、加一后的版本及该版本的完整展示快照共同提交。发送和重发都读取这份已保存快照，沿用原版本与原内容。W2 接管第 8 代时读取 T1 的原进度记录，下一次推进接着原版本分配，领取换代保留这条业务版本序列。

设 a 是版本 1 的核查已开始，b 是版本 2 的核查已完成。b 带有展示所需的开始时刻、当前阶段与完成结论，能够独立恢复从开始到完成的整份展示。报告服务从来件取得 T1、原版本及完整内容，核对其来源和资源映射，再在同一资源行锁下读取已保存的代次水位、业务版本与内容。先拒绝小于水位的代次，其余来件把水位推进到当前水位与传入代次的较大值，再核对操作身份与业务版本。较大业务版本才覆盖内容，较小版本保存迟到观察并保留当前内容。相同版本且内容一致，恢复该版本的原业务结果，后续准入使用本次事务推进后的当前代次水位；相同版本却内容不同，则保存冲突并核查。b 先提交版本 2 后，a 的版本 1 就只留下观察，展示保持 b。

资源把水位、展示内容、业务版本、操作输入及其结果共同提交。新代次即使只带来较旧业务版本，也会推进代次水位，迟到内容则留作观察。安装新领取水位时保留当前业务版本，避免换代重新开放旧进度。上面的单行 UPDATE 只表达代次条件；用于这份进度展示时，业务版本检查也要进入同一锁定事务或条件更新。若输入是必须逐步执行的增量，接收端则检查合法前驱，前驱尚缺时持久保存待处理输入，补齐后继续。第十章将沿这类增量展开缺口恢复。

为单独看清同代次内的倒序，图~\ref{fig:ch08-progress-version} 另取报告资源水位已为 8、业务版本初始为 0。W2 在购买库依次保存版本 1 的 a 和版本 2 的 b，发送时都带第 8 代，随后 b 先到报告服务。

\CEightFigure{ch08-progress-version}{同一领取代次内以业务版本处理倒序，较旧进度只保存观察}{fig:ch08-progress-version}

回到付款主线，实际支付平台通常只接收自己的支付身份，购买服务便从原 C 读取固定映射，通过平台的身份仲裁恢复付款。W1 在第 7 代发出的 \code{pay-C1} 在途中，W2 的第 8 代发送先到平台并形成 X1，W1 的旧请求随后到达。在第七章规定的保护范围内，两次到达都对应原平台身份和原输入，最终只有一条 3000 分流水，余额为 7000。

若 W2 把键改成 \code{pay-C1-epoch8}，平台会把它和旧 \code{pay-C1} 分别受理，余额可能变成 4000。本地 T1 始终只有一项，第 8 代本地写入也全部正确，重复付款已经发生。任务代次随接管改变，平台支付身份在 C1 受理后保持固定。

发送前检查租期可以减少失去资格的工作者发起新调用。已经通过检查、在途或已被平台受理的请求继续由平台处理。即使平台接收 token，高水位也是在观察到新代次以后排除旧写。同一 C1 的资金动作通过固定支付身份与输入绑定合流，旧请求和新请求共同恢复这一笔付款。

接管还会读出原 \code{provider\_retry\_before}，超过第七章的安全发送条件后转入查询或核查。续期和接管保留这个原截止。平台若在执行端检查固定截止 100，101 到达的旧尝试可被拒绝，恢复者继续查询原 C1 是否已经在此前形成 X1。租期到期依据任务时刻，调用过期依据尝试截止。固定身份最终拒绝则要求确认既往无成功，并排除所有已受理、在途及未来尝试成功。
